\documentclass{article}
\usepackage{amsmath}
\usepackage{amssymb}
\setlength{\parindent}{0pt}

\title{series}
\author{alexander}
\date{\today}
\begin{document}
\maketitle

a sequence we will say is an ordered list of terms, which we can denote by $a_1,a_2,a_3,\ldots$. a series is formed by adding the terms of a sequence: $a_1+a_2+a_3+\ldots$. for a finite sequence, this is simply a finite sum. for an infinite sequence, however, we cannot simply perform infinitely many additions and assume that the result is a finite number. instead, we consider the sequence of partial sums and ask whether those partial sums approach a finite limit. arithmetic and geometric sequences are particular types of sequences with simple relationships between consecutive terms.\\

in an arithmetic sequence, each term is obtained from the previous term by adding a constant difference, $a_{n+1}=a_n+d$, so the $n$th term can be written as $a_n=a+(n-1)d$. in a geometric sequence, each term is obtained from the previous term by multiplying by a constant ratio, $a_{n+1}=ra_n$, so the $n$th term can be written as $a_n=ar^{n-1}$. from these terms we can form the corresponding finite series, $\sum_{k=1}^{n}(a+(k-1)d)$ and $\sum_{k=1}^{n}ar^{k-1}$. there are other sequences that are neither arithmetic nor geometric; for now, however, we begin with these.\\

for a finite arithmetic series, the nth partial sum is:
\begin{enumerate}
	\item $S_n = a + (a + d) + (a + 2d) + \ldots + (a + (n - 1)d)$
	\item $S_n = (a + (n - 1)d) + (a + (n - 2)d) + \ldots + a$
	\item $2S_n = [a + (a + (n - 1)d)] + [(a + d) + (a + (n - 2)d)] + \ldots$
	\item $2S_n = n(2a + (n - 1)d)$
	\item $S_n = \frac{n}{2}(2a + (n - 1)d)$
\end{enumerate}

\newpage

for a finite geometric series, the nth partial sum is:
\begin{enumerate}
	\item $S_n - rS_n = (a + ar + ar^2 + \ldots + ar^{n-1}) - (ar + ar^2 + ar^3 + \ldots + ar^n)$
	\item $S_n - rS_n = a - ar^n$
	\item $S_n(1-r) = a-ar^n$
	\item $S_n = \frac{a-ar^n}{1-r} = \frac{a(1-r^n)}{1-r}$, $r\neq1$
\end{enumerate}

the restriction $r\neq1$ comes from the algebraic expression $\frac{a(1-r^n)}{1-r}$, since when $r=1$ the denominator is zero. however, the original geometric sum is still well-defined when $r=1$. in that case every term is $a$, so $S_n=na$. therefore, the restriction applies to this particular form of the formula, not to the original sum itself.\\

let us now consider the infinite series
\[
1+\frac{1}{2}+\frac{1}{4}+\frac{1}{8}+\frac{1}{16}+\frac{1}{32}+\ldots
\]
which can be written as
\[
\sum_{n=1}^{\infty}\frac{1}{2^{n-1}}.
\]

first, let us look at the series numerically by considering its partial sums:

\[
S_1=1=1.0
\]

\[
S_2=1+\frac{1}{2}=\frac{3}{2}=1.5
\]

\[
S_3=1+\frac{1}{2}+\frac{1}{4}=\frac{7}{4}=1.75
\]

\[
S_4=1+\frac{1}{2}+\frac{1}{4}+\frac{1}{8}=\frac{15}{8}=1.875
\]

\[
S_5=1+\frac{1}{2}+\frac{1}{4}+\frac{1}{8}+\frac{1}{16}
=\frac{31}{16}=1.9375
\]

\[
S_6=1+\frac{1}{2}+\frac{1}{4}+\frac{1}{8}+\frac{1}{16}+\frac{1}{32}
=\frac{63}{32}=1.96875
\]

the sequence of partial sums is therefore
\[
1,\frac{3}{2},\frac{7}{4},\frac{15}{8},\frac{31}{16},\frac{63}{32},\ldots
\]
or, numerically,
\[
1.0,1.5,1.75,1.875,1.9375,1.96875,\ldots
\]

we can observe that the terms of the series approach zero as $n$ increases. the partial sums, however, appear to approach the finite number $2$.\\

we can now use the formula for the finite geometric series to find an expression for the $n$th partial sum. in our example, the first term is $a=1$ and the common ratio is $r=\frac{1}{2}$. since $r\neq1$, we can use the formula:
\[
S_n=\frac{a(1-r^n)}{1-r}.
\]

substituting $a=1$ and $r=\frac{1}{2}$ gives
\[
S_n
=
\frac{1\left(1-\left(\frac{1}{2}\right)^n\right)}
{1-\frac{1}{2}}.
\]

since multiplying by $1$ does not change the numerator,
\[
S_n
=
\frac{1-\left(\frac{1}{2}\right)^n}
{1-\frac{1}{2}}.
\]

simplifying the denominator,
\[
1-\frac{1}{2}
=
\frac{2}{2}-\frac{1}{2}
=
\frac{1}{2}.
\]

therefore,
\[
S_n
=
\frac{1-\left(\frac{1}{2}\right)^n}
{\frac{1}{2}}
=
2\left(1-\left(\frac{1}{2}\right)^n\right).
\]

distributing the $2$ gives
\[
S_n
=
2-2\left(\frac{1}{2}\right)^n.
\]

we can simplify the second term:
\[
2\left(\frac{1}{2}\right)^n
=
\frac{2}{2^n}
=
\frac{2^1}{2^n}
=
2^{1-n}
=
\frac{1}{2^{n-1}}.
\]

therefore,
\[
\boxed{S_n=2-\frac{1}{2^{n-1}}}.
\]

we can now take the limit of the sequence of partial sums:
\[
\lim_{n\to\infty}S_n
=
\lim_{n\to\infty}
\left(
2-\frac{1}{2^{n-1}}
\right).
\]

as $n$ approaches infinity, $2^{n-1}$ becomes arbitrarily large, so
\[
\lim_{n\to\infty}\frac{1}{2^{n-1}}=0.
\]

therefore,
\[
\lim_{n\to\infty}S_n
=
2-0
=
2.
\]

thus, the sequence of partial sums approaches the finite number $2$, so this infinite geometric series converges to $2$.\\\\

the harmonic series is defined as
\[
\sum_{n=1}^{\infty}\frac{1}{n}
=
1+\frac{1}{2}+\frac{1}{3}+\frac{1}{4}+\frac{1}{5}+\ldots
\]

notice that the denominators of the harmonic series are the natural numbers:
\[
1,2,3,4,5,6,\ldots
\]
so the denominator grows linearly as $n$ increases.

this is different from the geometric series we just considered, whose denominators are powers of $2$:
\[
1,2,4,8,16,32,\ldots
\]
and therefore grow exponentially.

we can also notice that the geometric series is, in a sense, contained within the harmonic series. the harmonic series contains
\[
1,\frac{1}{2},\frac{1}{3},\frac{1}{4},\frac{1}{5},\frac{1}{6},\ldots
\]
while the geometric series contains
\[
1,\frac{1}{2},\frac{1}{4},\frac{1}{8},\frac{1}{16},\frac{1}{32},\ldots
\]

after the first two terms, each corresponding term of the harmonic series is larger than the corresponding term of the geometric series:
\[
\frac{1}{3}>\frac{1}{4},
\qquad
\frac{1}{4}>\frac{1}{8},
\qquad
\frac{1}{5}>\frac{1}{16},
\qquad
\frac{1}{6}>\frac{1}{32},
\ldots
\]

therefore, although both series have terms that approach zero, the terms of the harmonic series approach zero much more slowly. this suggests that simply approaching zero may not be enough for a series to converge; the rate at which the terms approach zero may also matter.\\

we can now make this observation more precise by grouping the terms of the harmonic series. after the first two terms, we group the remaining terms into groups containing $2,4,8,\ldots$ terms:
\[
1+\frac{1}{2}
+
\left(\frac{1}{3}+\frac{1}{4}\right)
+
\left(\frac{1}{5}+\frac{1}{6}+\frac{1}{7}+\frac{1}{8}\right)
+
\left(\frac{1}{9}+\cdots+\frac{1}{16}\right)
+\ldots
\]

consider the first group:
\[
\frac{1}{3}+\frac{1}{4}
>
\frac{1}{4}+\frac{1}{4}
=
\frac{2}{4}
=
\frac{1}{2}.
\]

the next group contains $4$ terms, and each term is at least $\frac{1}{8}$:
\[
\frac{1}{5}+\frac{1}{6}+\frac{1}{7}+\frac{1}{8}
>
\frac{1}{8}+\frac{1}{8}+\frac{1}{8}+\frac{1}{8}
=
\frac{4}{8}
=
\frac{1}{2}.
\]

the next group contains $8$ terms, and each term is at least $\frac{1}{16}$:
\[
\frac{1}{9}+\frac{1}{10}+\cdots+\frac{1}{16}
>
8\left(\frac{1}{16}\right)
=
\frac{8}{16}
=
\frac{1}{2}.
\]

we can therefore see a pattern: the number of terms in each group doubles, while the smallest term in each group is cut in half. the product of these two changes remains $\frac{1}{2}$, so each group contributes at least $\frac{1}{2}$ to the sum.\\

since each group contributes at least $\frac{1}{2}$, after adding $N$ such groups their total contribution is at least
\[
\underbrace{\frac{1}{2}+\frac{1}{2}+\cdots+\frac{1}{2}}_{N\text{ groups}}
=
\frac{N}{2}.
\]

as the number of groups increases without bound,
\[
\lim_{N\to\infty}\frac{N}{2}=\infty.
\]

therefore, the partial sums of the harmonic series can become arbitrarily large. they cannot approach a finite number, so the harmonic series diverges:
\[
\boxed{\sum_{n=1}^{\infty}\frac{1}{n}\text{ diverges}.}
\]

this gives us an important distinction between the two series. in the geometric series, the terms approach zero quickly enough that the partial sums approach the finite number $2$. in the harmonic series, the terms also approach zero, but they do so slowly enough that the partial sums continue to grow without bound.\\\\

for sufficiently large n the hierachy is roughly: constant < log < linear < exponential < combinatorial\\

p-series then comparison test

\end{document}
